\chapter{Projection and Retained Spaces}
\index{projection!orthogonal|textbf}
\index{retained subspace|textbf}
\index{compressed operator|textbf}

\label{ch:training-projection-retained-spaces}

\chapterstatus{pedagogical foundation for finite retained-space constructions.
The examples are illustrative mathematics and synthetic numerical checks, not
validated reductions of silicon.}

\section{Learning goals}

After working through this chapter, a reader should be able to

\begin{enumerate}
  \item construct an orthogonal projector from an orthonormal frame;
  \item distinguish a retained subspace from a basis used to represent it;
  \item define a compressed operator on the retained space;
  \item interpret an isolated Bloch band as a fiberwise retained space without
        confusing its scalar dispersion with a gauge or localized orbital;
  \item distinguish exact retention from an approximate effective-model class;
        and
  \item separate projection, disentanglement, basis transformation,
        localization, and truncation.
\end{enumerate}

The preceding chapters supplied finite representations, Hermitian eigensystems
and spectral projectors, quantum operators, periodic geometry, and a concrete
discretization. We can now state precisely what it means to retain only part of
a state space. Chapter~\ref{ch:training-hermitian-eigensystems} treated an
exact eigenspace of one matrix; this chapter adds selection rules, compression,
and fiberwise retained families.

\section{Orthogonal projection}

\begin{definition}[Orthogonal projection]
Let \(\mathcal S\subseteq\mathcal H\) be a closed subspace. Its orthogonal
projector \(\hat P_{\mathcal S}\) is the operator satisfying
\begin{equation}
\begin{aligned}
  \hat P_{\mathcal S}^{2}&=\hat P_{\mathcal S},
  &&\text{(idempotence)},\\
  \hat P_{\mathcal S}^{\dagger}&=\hat P_{\mathcal S},
  &&\text{(self-adjointness)},\\
  \operatorname{Ran}\hat P_{\mathcal S}&=\mathcal S,
  &&\text{(range condition)}.
\end{aligned}
\end{equation}
\end{definition}

For a finite orthonormal basis
\(\{\lvert\phi_a\rangle\}_{a=1}^{N}\) of \(\mathcal S\),
\begin{equation}
  \hat P_{\mathcal S}
  =\sum_{a=1}^{N}\lvert\phi_a\rangle\langle\phi_a\rvert.
\end{equation}
Each term retains the component along \(\lvert\phi_a\rangle\), while the sum
removes every component orthogonal to \(\mathcal S\).

If the orthonormal frame is stored as columns of
\(V\in\mathbb C^{D\times N}\), then
\begin{equation}
  V^{\dagger}V=I_N,
  \qquad
  P=VV^{\dagger}.
\end{equation}
Replacing \(V\) by \(VG\) for a unitary \(G\in U(N)\) changes the frame inside
\(\mathcal S\) but not the projector:
\begin{equation}
  (VG)(VG)^{\dagger}=VV^{\dagger}.
\end{equation}
The subspace and a selected basis of that subspace are therefore different
records.

\section{Compression to a retained space}
\index{compression!operator}

Let \(\widehat H\) act on its stated domain in a parent state space
\(\mathcal H\). When the required domain conditions hold, the compressed
operator on \(\mathcal S\) is
\begin{equation}
  \widehat H_{\mathcal S}
  =\left.\hat P_{\mathcal S}\widehat H\hat P_{\mathcal S}\right|_{\mathcal S}
  :\mathcal S\longrightarrow\mathcal S.
\end{equation}
In the frame \(V\), its finite matrix is
\begin{equation}
  H_{\mathcal S}=V^{\dagger}HV
\end{equation}
when \(H\) is a compatible ambient finite representation.

Selecting \(\mathcal S\) is a model or representation decision. Writing the
coordinates of \(\widehat H_{\mathcal S}\) in a basis is a separate operation.
Neither operation is synonymous with Wannier localization or with truncation of
matrix entries.

Consistently compressing an additive decomposition gives
\begin{equation}
  \hat P_{\mathcal S}\hat H\hat P_{\mathcal S}
  =\hat P_{\mathcal S}\hat T\hat P_{\mathcal S}
  +\hat P_{\mathcal S}\hat V\hat P_{\mathcal S}.
\end{equation}
In general, one may not compress \(\hat H\) while leaving \(\hat T\) in an
uncompressed parent representation. Projection may also change mathematical
form: a potential local in the original coordinate representation need not be
representable by multiplication by a scalar function on the retained space.

\section{Spectral retention}
\index{spectral reconstruction}

Assume for this pedagogical case that the relevant spectrum is discrete and
that the eigenstates form a complete orthonormal basis:
\begin{equation}
  \hat H\lvert\psi_i\rangle=E_i\lvert\psi_i\rangle,
  \qquad
  \hat H=\sum_i E_i\lvert\psi_i\rangle\langle\psi_i\rvert.
\end{equation}
For an energy cutoff \(E_{\mathrm{cut}}\), define
\begin{equation}
  \mathcal H^{(P)}
  =\operatorname{span}
   \left\{\lvert\psi_i\rangle\;\middle|\;E_i\leq E_{\mathrm{cut}}\right\},
  \qquad
  \hat P=\sum_{E_i\leq E_{\mathrm{cut}}}
  \lvert\psi_i\rangle\langle\psi_i\rvert.
\end{equation}
The retained Hamiltonian is
\begin{equation}
  \hat H^{(P)}
  =\left.\hat P\hat H\hat P\right|_{\mathcal H^{(P)}}
  =\sum_{E_i\leq E_{\mathrm{cut}}}
  E_i\lvert\psi_i\rangle\langle\psi_i\rvert.
\end{equation}
It exactly reproduces every retained eigenpair but does not specify the action
of \(\hat H\) outside \(\mathcal H^{(P)}\). It is an exact restriction of the
parent Hamiltonian, not yet an approximate effective model.

This construction is pedagogical. A periodic crystal uses Bloch-fiber spaces,
and entangled bands may require construction of a smooth retained family rather
than selection by one global energy cutoff.

\section{Isolated bands as fiberwise retained spaces}
\index{isolated band!retained space}
\index{band gap!sampled isolation}

For a periodic parent, retention is performed over a family of Bloch fibers.
Let
\begin{equation}
  \widehat H(k)\lvert u_{n k}\rangle
  =E_n(k)\lvert u_{n k}\rangle.
\end{equation}
A nondegenerate band is isolated over a reciprocal domain when its eigenvalue
remains separated from the excluded spectrum throughout that domain. The
associated rank-one projector is
\begin{equation}
  \widehat P_n(k)
  =\lvert u_{n k}\rangle\langle u_{n k}\rvert.
  \label{eq:training-isolated-band-projector}
\end{equation}
Changing the phase by
$\lvert u_{n k}\rangle\mapsto e^{i\theta(k)}\lvert u_{n k}\rangle$ changes the
frame but leaves $\widehat P_n(k)$ unchanged.

The retained Hamiltonian in this one-dimensional fiber space is the scalar
matrix
\begin{equation}
  H_n^{(P)}(k)
  =\langle u_{n k}\rvert\widehat H(k)\lvert u_{n k}\rangle
  =E_n(k).
  \label{eq:training-isolated-band-scalar-operator}
\end{equation}
This explains both the usefulness and the limitation of scalar isolated-band
reduction. The dispersion $E_n(k)$ completely specifies the Hamiltonian's
action inside each rank-one eigenspace. It does not specify a phase-connected
family of eigenvectors, a Wannier function, a localization center, or a
real-space probability density.

On a finite reciprocal mesh, a positive minimum sampled adjacent gap is a
conditioning diagnostic. It does not prove that the gap remains positive
between sample points. Near a degeneracy, energy ordering can exchange state
identity even when the ordered eigenvalue table remains easy to compute. A
composite retained projector may then remain meaningful after the individual
scalar labels cease to be stable.

Chapter~\ref{ch:training-isolated-band-fourier-reduction} starts from the
scalar retained operator in
Equation~\eqref{eq:training-isolated-band-scalar-operator} and derives its
finite hopping representation. Chapter~\ref{ch:training-composite-band-gauge-alignment}
then retains a higher-rank frame and makes gauge transport, alignment, the
Bloch--Wannier bridge, and block locality explicit. Neither Fourier
transformation is an additional projection or, by itself, a localization
calculation.

\section{Exact retained operator versus effective model}
\index{effective Hamiltonian!retained-operator distinction}

Reserve \(\hat H_{\mathrm{eff}}(\boldsymbol\theta)\) for an operator belonging
to a prescribed effective-model class with parameters \(\boldsymbol\theta\).
Then
\begin{equation}
  \underbrace{\hat H^{(P)}}_{\text{exact retained operator}}
  \qquad\text{and}\qquad
  \underbrace{\hat H_{\mathrm{eff}}(\boldsymbol\theta)}_{
    \text{approximate model on an identified retained space}}
\end{equation}
are different constructions.

The first is fixed after the parent, retained space, and compression rule are
fixed. The second additionally selects a model class and usually a fitting or
reduction criterion. Agreement of selected eigenvalues does not by itself show
that the effective operator reproduces the retained operator in the coordinates
needed by later perturbations.

\section{Operations that must not be collapsed}
\index{disentanglement!distinction}
\index{truncation!distinction}

\begin{description}
  \item[Projection.] Applies an idempotent map onto a declared subspace.
  \item[Spectral selection.] Chooses eigenvectors according to a declared
        spectral rule in a setting where that rule is meaningful.
  \item[Disentanglement.] Constructs a retained subspace from a larger window
        when target bands are not globally isolated.
  \item[Basis transformation.] Changes coordinates inside an already identified
        space without changing the represented abstract operator.
  \item[Localization.] Selects or optimizes a basis or gauge according to a
        locality criterion; it is not automatically a reduction in rank.
  \item[Truncation.] Removes matrix elements, basis functions, hopping blocks,
        or other represented content according to a separate rule.
  \item[Model fitting.] Selects parameters in a declared admissible class using
        an objective and data split.
\end{description}

Several operations may appear in one workflow, but each introduces different
assumptions and error sources. Calling all of them ``projection'' prevents those
errors from being evaluated separately.

\section{A finite projector calculation}

\begin{pythonexample}{compressing a synthetic Hermitian operator}
import numpy as np

# Define a synthetic ambient Hermitian operator on a four-coordinate space.
H = np.array(
    [
        [0.0, 0.2, 0.0, 0.0],
        [0.2, 1.0, 0.1, 0.0],
        [0.0, 0.1, 2.0, 0.3],
        [0.0, 0.0, 0.3, 3.0],
    ],
    dtype=np.complex128,
)

# Use two orthonormal columns as the retained frame.
V = np.eye(4, dtype=np.complex128)[:, :2]
P = V @ V.conjugate().T
H_retained = V.conjugate().T @ H @ V

# Check projector identities and the compressed representation explicitly.
np.testing.assert_allclose(P @ P, P)
np.testing.assert_allclose(P.conjugate().T, P)
np.testing.assert_allclose(H_retained, V.conjugate().T @ P @ H @ P @ V)
\end{pythonexample}

The example verifies exact identities for synthetic arrays. It does not justify
this retained frame for a physical problem.

\section{Questions to record before retention}

A retained-space calculation should state

\begin{enumerate}
  \item the parent state space and operator;
  \item the rule that defines the retained space;
  \item the retained rank and any conditioning requirement;
  \item the projector or frame, including basis and gauge conventions;
  \item the treatment of degeneracies or entangled bands;
  \item the operator terms compressed through the same map;
  \item the relation between the exact retained operator and any fitted model;
        and
  \item the evidence and tolerance used to assess each approximation.
\end{enumerate}

\section{Exercises}

\begin{enumerate}
  \item Prove that \(P=VV^{\dagger}\) is an orthogonal projector when
        \(V^{\dagger}V=I\).
  \item Show that \(V\) and \(VG\) define the same projector for unitary \(G\),
        and derive the corresponding covariance of \(V^{\dagger}HV\).
  \item Construct a local diagonal matrix \(V\) in an ambient basis and show
        that its compression can be dense in another retained frame.
  \item Explain why deleting small entries from \(H\) is not the same operation
        as projecting onto a retained subspace.
  \item State why a spectral cutoff can become ambiguous near a degeneracy or
        for entangled bands.
\end{enumerate}

\section{Evidence boundary}

Projector identities and exact recovery of retained eigenpairs can establish
mathematical or software properties of a declared construction. They do not
show that the retained space preserves the structures needed by a later
continuum model or that a fitted effective Hamiltonian is scientifically
adequate. Chapter~\ref{ch:training-evidence-vocabulary} supplies the vocabulary
needed to state those stronger claims without conflation.
